% ============================================================
%  АНОТАЦІЯ УКРАЇНСЬКОЮ ТА АНГЛІЙСЬКОЮ МОВАМИ, СПИСОК ПУБЛІКАЦІЙ
%  Жанрова вимога ДСТУ: розгорнута анотація (2--3 с.) укр. і англ.,
%  ключові слова, перелік публікацій здобувача за темою дисертації.
%  Цей файл — структурований каркас; прозу заповнюємо після розділів.
% ============================================================

\chapter*{Анотація}
\addcontentsline{toc}{chapter}{Анотація}

\noindent\textit{Мельниченко О. <Прізвище, ініціали по батькові>.} \thesistitle{} -- Кваліфікаційна наукова праця на правах рукопису.

\medskip

\noindent Дисертація на здобуття наукового ступеня доктора технічних наук за спеціальністю <шифр і назва спеціальності, напр.\ 122 -- Комп'ютерні науки / 125 -- Кібербезпека та захист інформації>. -- <Назва закладу>, Київ, 2026.

\bigskip

% --- Розгорнута анотація (каркас, ~2--3 с.) ---
% Блоки, які треба розкрити прозою:
%  1) Актуальність і наукова проблема.
%  2) Мета й завдання.
%  3) Об'єкт і предмет.
%  4) Стислий виклад наукової новизни (8 пунктів).
%  5) Практичне значення та впровадження.
%  6) Структура й обсяг роботи.
У дисертації розв'язано науково-прикладну проблему приватного контекстного захисту учасників цифрової комунікації від ризикових інформаційних впливів в умовах наскрізного шифрування та постквантового переходу. \emph{[Розгорнути: суть проблеми, основні результати, новизна, практичне значення -- після написання розділів.]}

\bigskip

\noindent\textbf{Ключові слова:} захист контекстної приватності, оцінювання інформаційного ризику, контекстна безпека систем обміну повідомленнями, наскрізне шифрування, постквантова криптографія, OPAQUE, асиметричний парольно-автентифікований обмін ключами, ML-KEM-768, серверно-сліпий канал нагляду, захисні рішення з урахуванням політик, відтворюваність, інформаційна технологія.

% ------------------------------------------------------------

\chapter*{Abstract}
\addcontentsline{toc}{chapter}{Abstract}

\begin{english}
\noindent\textit{Melnychenko O.} Methodology and information technology for protecting a person's contextual privacy under inferential threats of the digital environment. -- Qualifying scientific work on the rights of a manuscript.

\medskip

\noindent Thesis for the degree of Doctor of Technical Sciences. -- <Institution>, Kyiv, 2026.

\bigskip

\noindent The dissertation addresses the scientific and applied problem of private, context-aware protection of communication participants against harmful informational influence under end-to-end encryption and the post-quantum transition. \emph{[To be expanded after the chapters are written.]}

\bigskip

\noindent\textbf{Keywords:} private contextual protection, informational risk assessment, messenger-native safety, end-to-end encryption, post-quantum cryptography, OPAQUE, augmented PAKE, ML-KEM-768, server-blind supervision relay, policy-aware safety decisions, reproducibility, information technology.
\end{english}

% ------------------------------------------------------------

\section*{Список публікацій здобувача за темою дисертації}
\addcontentsline{toc}{section}{Список публікацій здобувача за темою дисертації}

% TODO(author): впорядкувати за ДСТУ 8302; розділити на групи:
%  -- статті у наукових фахових виданнях України та виданнях, проіндексованих у Scopus/WoS;
%  -- праці апробаційного характеру (тези, доповіді);
%  -- авторські свідоцтва / реєстрації ПЗ за наявності.
\begin{enumerate}
  \item Melnychenko O. \emph{Hybrid PQ-OPAQUE: A Three-Message Augmented PAKE with Ephemeral ML-KEM-768} // <Journal of Information Security and Applications (подано/прийнято)>, 2026. \emph{[уточнити статус, том, сторінки, DOI]}
  \item \emph{[додати решту публікацій за темою дисертації]}
\end{enumerate}
